% v. 3.4 2022/03/03 by sl:
%
% - changed DOI numeration of articles
% - changed default values of page,year,nr,volume,issue
%
% v. 3.3 2021/08/20 by sl:
%
% - added macros for distinguishing arxiv/ios final versions 
%
% v. 3.1 2021/01/05 by sl:
%
% - updated info about the last version - just before sect. 1.1
% - macro \year replaced by \publyear
% - tested with tikz - section 2.6 extended with a tikz example
%
%
%
% v. 3.0 2020/12/11 by sl:
%
% - updated info about the last version
% - new subsubsection 2.1.1 concerning the article meta-data
% - updated the bibliography section
%
%
%
%$Id: figuide.tex,v 2.0 2016/11/04 by szczuka@mimuw.edu.pl$


\documentclass{fundam}


%%%%%%%  parameters to be filled in by copy-editor  %%%%%%%%%%

\setcounter{page}{1}
\publyear{22}
\papernumber{2102}
\volume{185}
\issue{1}

%\finalVersionForARXIV
%\finalVersionForIOS

%%%%%%%%%%%%%%%%%%%%%%%%%%%%%%%%%%%%%%




\usepackage{url} % takes care of hyperlinks, preferred over hyperref
\usepackage[ruled,lined]{algorithm2e}% provides Algorithm environment
\usepackage{graphicx}% allows for inclusion of graphic files (figures)
\usepackage{tikz}
\usetikzlibrary{automata, positioning, arrows}

\usepackage{amsfonts}
\usepackage{amssymb}
\usepackage{framed}
%\usepackage{color,soul}
\usepackage{listings}
%%%
\usepackage{makecell}
\usepackage{ctable}
\usepackage{amsmath}
%%%
\usepackage[
n,
operators,
advantage,
sets,
adversary,
landau,
probability,
notions,	
logic,
ff,
mm,
primitives,
events,
complexity,
asymptotics,
keys]{cryptocode}


\usepackage{mathrsfs}

%\newcommand{\todo}[1]{{\color{red}\bf TODO:{ #1}!!!}}
\newcommand{\vr}[1]{\mathbf{#1}}
\newcommand{\Lat}{\mathcal{L}}

\newcommand{\Exp}[2]{\mathbf{Exp}_{#1}^{#2}}
\newcommand{\Next}{\mathcal{O}.\mathsf{Next}}
\newcommand{\Upd}{\mathcal{O}.\mathsf{Upd}}
\newcommand{\UpdC}{\mathcal{O}.\mathsf{Upd}\widetilde{\mathsf{C}}}
\newcommand{\Corr}{\mathcal{O}.\mathsf{Corr}}
\newcommand{\Chall}{\mathcal{O}.\mathsf{Chall}}
\newcommand{\OEnc}{\mathcal{O}.\mathsf{Enc}}
\newcommand{\ODec}{\mathcal{O}.\mathsf{Dec}}
\newcommand{\Set}[1]{\mathcal{#1}}
\newcommand{\opuni}{\mathsf{op}\md\mathsf{uni}}
\newcommand{\funi}{\mathsf{f}\md\mathsf{uni}}
\newcommand{\uekg}{\mathsf{UE.KG}}
\newcommand{\uetg}{\mathsf{UE.TG}}
\newcommand{\ueenc}{\mathsf{UE.Enc}}
\newcommand{\uedec}{\mathsf{UE.Dec}}
\newcommand{\ueupd}{\mathsf{UE.Upd}}
\newcommand{\UE}{\mathsf{UE}}
\newcommand{\EUmod}[1]{\mathsf{rand}\md\mathsf{ind}\md\mathsf{eu}\md\mathsf{cpa}\md{#1}}
\newcommand{\EUS}{\mathsf{rand}\md\mathsf{ind}\md\mathsf{eu}\md\mathsf{cpa}}
\renewcommand{\indcpa}{\mathsf{ind}\md\mathsf{cpa}}
\newcommand{\Adv}{\mathsf{Adv}}

\newcommand{\params}{\mathsf{parameters}}

\newcommand{\md}{{\text{-}}}

\newcommand{\ct}{\mathsf{ct}}
\newcommand{\qid}{\mathsf{qid}}
\newcommand{\m}{\mathbf{m}}
\newcommand{\epoch}{\mathsf{e}}
\newcommand{\tok}{\Delta}
\newcommand{\tx}[1]{\mathsf{#1}}
\renewcommand{\opuni}{\mathsf{b}\md\mathsf{uni}}

\renewcommand{\Lat}{\mathcal{L}}
\newcommand{\Fun}{\mathcal{F}}
\newcommand{\Par}{\mathcal{P}}
\newcommand{\vol}{\text{vol}}

\newcommand{\sivp}{\ensuremath{\mathsf{SIVP}}}
\newcommand{\Distr}{\mathcal{D}}
\newcommand{\ii}{\mathcal{U}(n)}

\newtheorem{concl}[definition]{Conclusion}
\renewcommand{\sample}{\overset{\$}{\gets}}

\newcommand\quotient[2]{
	\mathchoice
	{% \displaystyle
		\hspace{0.1cm}\text{\raise1ex\hbox{$#1$}\Big/\lower1ex\hbox{$#2$}}
	}
	{% \textstyle
		#1\,/\,#2
	}
	{% \scriptstyle
		#1\,/\,#2
	}
	{% \scriptscriptstyle  
		#1\,/\,#2
	}
}
%%%%%%%%%%%%%%%%%%%%%%%%%%%%%%%%%%%%%%%%%%%%%%%%%%%%%%%%%%%%%%%%%%%%%%%%%%%%%%


\newcommand{\todo}[1]{{\color{red}\textbf{TODO\{#1!!!\}}}}

\begin{document}



\title{Unidirectional Key Update in Updatable Encryption, Revisited}

\address{Military University of Technology, Kaliskiego 2 00-908 Warsaw}

\author{Mariusz Jurkiewicz, Kamila Prabucka\thanks{Supported by the Grant DOB/002/RON/ID1/2018 by The National Centre for Research and Development.}\\
	Military University of Technology, Warsaw, Poland\\
	\{mariusz.jurkiewicz, kamila.prabucka\}@wat.edu.pl}

\maketitle

\runninghead{M. Jurkiewicz, K. Prabucka}{Unidirectional Key Update in Updatable Encryption, RevisitedI}

\begin{abstract}
	In this paper we construct a new efficient updatable encryption ($\UE$) scheme based on $\tx{FrodoPKE}$ learning with errors key encapsulation \cite{Frodo}. We analyse the security of the proposed scheme in the backward-leak uni-directional setting  within the $\mathsf{rand}\md\mathsf{ind}\md\mathsf{eu}\md\mathsf{cpa}$ model. Since the underlying computationally hard problem here is LWE, the scheme is secure against both classical and quantum attacks. 
	% In addition, we analyze the existing security requirement and provide the new one that turns out to be strictly stronger than its more general counterpart. 
\end{abstract}

\begin{keywords}
	Updatable Encryption; Unidirectional Key Update;  Lattices; LWE problem.
\end{keywords}


\section{Introduction}
Updatable encryption ($\UE$) allows a third party to periodically rotate encryption keys from one epoch to another without the need to download, decrypt, re-encrypt and upload already encrypted data by a client. Updating those outsourced ciphertexts is carried out via the use of so-called update tokens which in turn are generated during key rotation and can be sent (publicly) to the third party.

\vspace{0.1cm}
In this paper we construct a new efficient $\UE$ scheme and consider its security in $\mathsf{rand}\md\mathsf{ind}\md\mathsf{eu}\md\mathsf{cpa}$ model. Since the underlying computationally hard problem here is LWE, the scheme is secure against both classical and quantum attacks.  In addition, we analyse the existing security requirement and provide the new one that turns out to be strictly stronger than its more general counterpart.  


\vspace{0.1cm}
An updatable encryption scheme for message space $\Set{M}$ consists of a tuple of PPT  algorithms $(\uekg,\uetg,\ueenc, \uedec,$ $\ueupd)$ that operate in epochs, the epoch starts at $0$. The \textit{key generation} algorithm $\uekg$ outputs an epoch key $\key_{\epoch}$. The \textit{token generation} algorithm $\uetg$ takes as input two epoch keys $\key_{\epoch}$, $\key_{\epoch+1}$ and outputs an update token $\tok_{\epoch+1}$, the update token can be used to move ciphertexts from epoch $\epoch$ to $\epoch+1$. The encryption algorithm $\ueenc$ takes as input an epoch key $\key_{\epoch}$ and a message $\m$ and outputs a ciphertext $\ct_{e}$. The decryption algorithm $\uedec$ takes as input an epoch key $\key_{\epoch}$ and a ciphertext $\ct_{\epoch}$ and outputs a message $\m'$. The update algorithm $\ueupd$ takes as input an update token $\tok_{\epoch+1}$ and aciphertext $\ct_{\epoch}$ from epoch $\epoch$ and outputs an updated ciphertext (see \cite{jiang2020} for more details).

\vspace{0.1cm}
We say that $\UE$ scheme is \textit{correct} if for any message $\m\in\Set{M}$ and any epochs $\epoch_{1}\leq \epoch_{1} \leq \epoch_{\max}$, we have
$$\Pr\left[  \uedec(\key_{\epoch_{2}}, \ct_{\epoch_{2}}) \neq \m \right]\leq \negl$$
where $\key_{\epoch_{1}}, \ldots \key_{\epoch_{2}} \sample \uekg(1^{n})$, $\ct_{\epoch_{1}}\sample \ueenc(\key_{\epoch_{1}}, \m)$; $\tok_{j}\gets \uetg(\key_{j-1}, \key_{j})$,\\ $\ct_{j}\gets \ueupd(\tok_{j}, \ct_{j-1})$, for $j\in \set{\epoch_{1}\tx{+}1, \ldots ,\epoch_{2} }$.
%The formal definition of UE is intuitive and so we refer to \cite{jiang2020} for details. On the other hand, 

\vspace{0.1cm}
The security notion(s) for $\UE$ is based on \textit{leakage sets} given by \cite{klooss2019r}, \cite{jiang2020}. These sets are used to track epochs in which the adversary corrupted a key, a token, or a version of challenge-ciphertext:
\begin{itemize}
\item $\Set{K}$: set of epochs in which the adversary corrupted the epoch key through $\Corr$; 
\item $\Set{T}$: set of epochs in which the adversary corrupted the update token through $\Corr$; 
\item $\Set{C}$: set of epochs in which the adversary obtained a challenge-equal ciphertext through $\Chall$ or $\UpdC$.
\end{itemize}

In addition, it is also the need to track ciphertexts and their updates that can be known to the adversary:
\begin{itemize}
\item $\Set{L}$: Set of non-challenge ciphertexts $(\qid, \ct, \epoch ; \m)$, where query identifier $\qid$ is a counter incremented by each invocation of $\OEnc$, $\ct$ is the given ciphertext, $\epoch$ is the epoch when the query is asked, and $\m$ is the underlying plaintexts. The adversary learned these ciphertexts from $\OEnc$ or $\Upd$.
\item $\widetilde{\Set{L}}$: Set of challenge-equal ciphertexts $(\widetilde{\ct}, \epoch)$, where $\widetilde{\ct}$ is the given challenge-equal ciphertext, and $\epoch$ is the epoch when the query is asked. The adversary learned these ciphertexts from $\Chall$ or $\UpdC$
\end{itemize} 


In \cite{jiang2020} the sets $\Set{L}$ and $\widetilde{\Set{L}}$, along with their extensions, are used only in the deterministic versions of $\UE$, and concerning the randomized UE there are defined two plaintexts-based sets $\Set{Q}^{*}$ and $\widetilde{\Set{Q}}^{*}$. 
%Due to this distinction Y. Jiang obtains a one-to-one correspondence between ciphertexts and the associated messages, independently of the case. 
This distinction lets the author of \cite{jiang2020} to obtain a one-to-one correspondence between ciphertexts and the associated messages, independently of the case. 
However, although in this paper we construct and consider a randomized $\UE$ scheme we show that such strict assumptions are not necessary and it is sufficient to only take advantage of the former group of these sets.


In addition, since the adversary is able to learn or infer information via its known tokens, there are also defined $\Set{K}^{*}$, $\Set{T}^{*}$, $\Set{C}^{*}$, $\Set{L}^{*}$ and $\widetilde{\Set{L}}^{*}$ as the extended sets of $\Set{K}$, $\Set{T}$, $\Set{C}$, $\Set{L}$ and $\widetilde{\Set{L}}$.

\vspace{0.1cm}
\noindent\textbf{Key Leakage.} According to \cite{jiang2020} the following three cases can be separately consider as for the key update:
\begin{itemize}
\item  the knowledge of a pair $(\key_{\epoch}$, $\tok_{\epoch+1})$ does not enable to generate $\key_{\epoch+1}$ ($\tx{no}$-directional key update);
\item the knowledge of a pair $(\key_{\epoch}$, $\tok_{\epoch+1})$ enables to generate $\key_{\epoch+1}$ ($\tx{uni}$-directional key update);
\item the knowledge of $(\key_{\epoch}$, $\tok_{\epoch+1})$ or $(\key_{\epoch+1}$, $\tok_{\epoch+1})$ enables to generate $\key_{\epoch+1}$ or $\key_{\epoch}$, resp. ($\tx{bi}$-direct. key upd.).
\end{itemize}


At first glance it appears that $\tx{bi}$-directional key update implies $\tx{uni}$-directional key update However, it is not actually the case as here the unidirectionality means precisely that update is possible just in only one direction. This means, in particular, that this property does not hold for $\tx{bi}$-directional key update.  
Ryo Nishimaki noticed in \cite{nishimaki2022direction} that $\tx{no}, \tx{uni}, \tx{bi}$-directional key updates do not cover all the possible cases and, furthermore, that the opposite variant of the unidirectionality, where  the knowledge of $(\key_{\epoch+1}$, $\tok_{\epoch+1})$ enables to obtain $\key_{\epoch}$, constitutes an important separate case. % and must be can be separately  
Consequently, the notion of $\tx{uni}$-directional key update has been split into \textit{forward-leak uni-directional key update} $\funi$ (the old $\tx{uni}$-, i.e. $\Set{K}_{\funi}^{*}=\Set{K}_{\mathsf{uni}}^{*}$) and \textit{backward-leak uni-directional key update} $\opuni$. 
If $l+1\in \ZZ_{>0}$ denotes the maximal number of epochs, the set $\Set{K}_{\opuni}^{*}$ is defined as follows


$$\Set{K}_{\opuni}^{*}\gets \set{\epoch\in [0,l]\mid \tx{CorrK}(\epoch)=\tx{true}},$$ $$ \text{ where } \tx{true}\gets\tx{CorrK}(\epoch)\Leftrightarrow (\epoch \in\Set{K}) \lor (\tx{CorrK(\epoch +1 )\land (\epoch +1 )\in \Set{T}}).$$
It turns out that $\opuni$-directional key update is stronger than $\tx{bi}$-directional key update in terms of security. Furthermore, Yao Jiang Galteland and Jiaxin Pan obtained in \cite{jiang2023backward} an astonishing result that security in the backward-leak uni-directional key updates setting is equivalent to the no-directional one.

\vspace{0.1cm}
\noindent\textbf{Token Leakage.} A token that is known to the adversary is either a corrupted token or a token inferred from two consecutive epoch keys, so $\Set{T}_{\opuni}^{*} = \Set{T}_{\opuni}^{*} (\Set{T}, \Set{K}_{\opuni}^{*}),$
$$\Set{T}_{\opuni}^{*} \gets \set{\epoch\in [0,l]\mid  (\epoch \in \Set{T}) \lor (\epoch \in \Set{K}_{\opuni}^{*} \land (\epoch -1)\in  \Set{K}_{\opuni}^{*})  }.$$

\vspace{0.1cm}
\noindent\textbf{Challenge-Equal Ciphertext Leakage.} The adversary learned all challenge-equal ciphertexts in epochs from $\Set{C}$. Furthermore, the adversary can infer challenge-equal ciphertexts via tokens. It is easily seen that a form of $\Set{C}^{*}$ depends on if ciphertexts are $\tx{uni}$- or $\tx{bi}$-directional updated, and it is independent of assumed $\opuni$-directional updated of keys (and tokens). Let $\tx{cc}\in \set{\tx{uni}, \tx{bi}}$, then the method of computing the sets $\Set{C}_{\opuni, \tx{cc}}^{*}$ is presented in \cite{jiang2020}.

\vspace{0.1cm}
\noindent\textbf{Oracles in security games for updatable encryption}. The security game for a variant of $\UE$ we are interested herein uses the following oracles $\mathbf{Setup}$, $\OEnc$, $\ODec$, $\Next$, $\Upd$, $\Corr$, $\Chall$ and $\UpdC$.  




\begin{pchstack}[center]
	
	\begin{pcvstack}%
		\procedure{$\mathbf{Setup}(1^{n})$}{%
			\pcln \key_{0}\sample \uekg(1^{n})\\
			\pcln \tok_{0}\gets \perp; \epoch, \qid, \tx{twf}\gets 0\\
			\pcln \Set{L}, \widetilde{\Set{L}}, \Set{C}, \Set{K}, \Set{T}\gets \emptyset }
		
		\pcvspace
		
		\procedure{$\OEnc(\m)$}{%
			\pcln  \qid \gets \qid +1   \\
			\pcln  \ct \sample \ueenc(\key_{\epoch}, \m)  \\
			\pcln  \Set{L} \gets \Set{L}\cup \set{\qid, \ct, \epoch; \m}  \\
			\pcln   \pcreturn \ct }
		
	\end{pcvstack}
	
	\pchspace
	\hspace{1.5cm}
	
	\begin{pcvstack}%
		
		
		\procedure{$\Corr(\tx{inp}, \hat{\epoch})$}{%
			\pcln  \pcif \hat{\epoch}>\epoch~\textbf{then}\\
			\pcln ~~~~~\pcreturn \perp\\
			\pcln\pcif \tx{inp}=\tx{key}~\textbf{then}\\
			\pcln ~~~~~\Set{K}\gets \Set{K}\cup \set{\hat{\epoch}}\\
			\pcln ~~~~~\pcreturn \key_{\hat{\epoch}}\\
			\pcln \pcif \tx{inp}=\tx{token}~\textbf{then}\\
			\pcln ~~~~~\Set{T}\gets \Set{T}\cup \set{\hat{\epoch}}\\
			\pcln ~~~~~\pcreturn \tok_{\hat{\epoch}}}	
		
	\end{pcvstack}
	
\end{pchstack}


%%%%%%

\begin{pchstack}[center]
	
	\begin{pcvstack}%
		
		\procedure{$\ODec(\ct)$}{%
			\pcln  \m' ~\tx{or}~\perp \gets \uedec(\key_{\epoch}, \ct) \\
			\pcln \pcif (\ct, \epoch)\in \widetilde{\Set{L}}^{*}~\textbf{then}\\
			\pcln ~~~~~\tx{twf}\gets 1\\
			\pcln   \pcreturn \m'~\textbf{or}~\perp }
		
		\pcvspace
		
		\procedure{$\Next()$}{%
			\pcln  \qid \gets \qid +1\\
			\pcln \key_{\epoch} \sample \uekg(1^{n})\\
			\pcln \tok_{\epoch}\gets \uetg(\key_{\epoch-1}, \key_{\epoch})\\
			\pcln\pcif \tx{phase}=1 ~\textbf{then}\\
			\pcln   ~~~~~\widetilde{\ct}_{\epoch} \gets \ueupd(\tok_{\epoch}, \widetilde{\ct}_{\epoch\md1})}
		
		\pcvspace
		
		\procedure{$\Upd(\ct_{\epoch-1})$}{%
			\pcln  \pcif (j, \ct_{\epoch-1}, \qid-1, \m)\not\in \Set{L}~\textbf{then}\\
			\pcln~~~~~\pcreturn\perp\\
			\pcln \ct_{\epoch}\gets \ueupd(\tok_{\epoch}, \ct_{\epoch-1})\\
			\pcln \Set{L}\gets \Set{L}\cup \set{(j, \ct_{\epoch}, \epoch; \m)}\\
			\pcln   \pcreturn \ct_{\epoch}}	
		
	\end{pcvstack}
	
	\pchspace
	\hspace{1.5cm}
	
	\begin{pcvstack}%
		
		\procedure{$\Chall(\bar{\m}, \bar{\ct}, \tx{b})$}{%
			\pcln\pcif  \tx{phase}=1 ~\textbf{then}  \\
			\pcln ~~~~~\textbf{return}~\perp  \\
			\pcln  \tx{phase}\gets 1; \widetilde{\epoch}\gets \epoch  \\
			\pcln\pcif  (\cdot, \bar{\ct}, \widetilde{\epoch}-1, \bar{\m}_{1})\not\in \Set{L} ~\textbf{then}  \\
			\pcln ~~~~~\textbf{return}~\perp  \\
			\pcln\pcif  \tx{b}=0 ~\textbf{then}\\
			\pcln ~~~~~\widetilde{\ct}_{\widetilde{\epoch}} \gets \ueenc(\key_{\widetilde{\epoch}}, \bar{\m}) \\
			\pcln\pcelse \\
			\pcln ~~~~~\widetilde{\ct}_{\widetilde{\epoch}} \gets \ueupd(\tok_{\widetilde{\epoch}}, \bar{\ct}) \\
			\pcln \Set{C}\gets \Set{C}\cup \set{\widetilde{\epoch}}\\
			\pcln \widetilde{\Set{L}} \gets \widetilde{\Set{L}} \cup \set{(\widetilde{\ct}_{\widetilde{\epoch}}, \widetilde{\epoch})}\\
			\pcln   \pcreturn \widetilde{\ct}_{\widetilde{\epoch}}}
		
		\pcvspace
		
		\procedure{$\UpdC()$}{%
			\pcln \pcif \tx{phase}\neq 1~\textbf{then}\\
			\pcln ~~~~~\pcreturn \perp\\
			\pcln \Set{C}\gets \Set{C}\cup\set{\epoch}\\
			\pcln \widetilde{\Set{L}} \gets \widetilde{\Set{L}} \cup \set{(\widetilde{\ct}_{\epoch}, \epoch)}\\
			\pcln   \pcreturn \widetilde{\ct}_{\epoch} }		
		
		
	\end{pcvstack}
	
\end{pchstack}
\vspace{0.1cm}
Although our $\UE$ scheme is of a general randomized version, it can be easily reformulated to a deterministic scheme. In spite of this fact, we decided to focused only on this general case and, therefore, adjust the security model to this specific case. The security itself is based on the classical semantic approach, where an adversary  has to distinguish between two given experiments, namely $\Exp{\UE,\adv}{\EUmod{0}}$ and $\Exp{\UE,\adv}{\EUmod{1}}$. Furthermore, it must be emphasized that the presented scheme turns out to be indistinguishable over $\tx{CPA}$ attack.
\begin{pchstack}[center]
	\procedure{$\Exp{\UE,\adv}{\EUmod{0}}$}{%
		\pcln\textbf{do} ~\textbf{Setup};~\tx{phase}\gets 0\\
		\pcln \tx{b}'\gets \adv^{\tx{oracles}(\tx{b}=0)}(1^{n})\\
		\pcln\pcif \Set{K}^{*}\cap \Set{C}^{*}\neq \emptyset~\textbf{then}\\
		\pcln~~~~~\tx{twf}\gets 1\\
		\pcln\pcif \tx{twf}= 1~\textbf{then}\\
		\pcln ~~~~~\tx{b}'\sample \bin\\
		\pcln   \pcreturn \tx{b}' }		
	
	\pchspace
	
	\hspace{2cm}
	
	\pchspace
	\procedure{$\Exp{\UE,\adv}{\EUmod{1}}$}{%
		\pcln \textbf{do} ~\textbf{Setup};~\tx{phase}\gets 0\\
		\pcln \tx{b}'\gets \adv^{\tx{oracles}(\tx{b}=1)}(1^{n})\\
		\pcln\pcif \Set{K}^{*}\cap \Set{C}^{*}\neq \emptyset~\textbf{then}\\
		\pcln~~~~~\tx{twf}\gets 1\\
		\pcln\pcif \tx{twf}= 1~\textbf{then}\\
		\pcln ~~~~~\tx{b}'\sample \bin\\
		\pcln   \pcreturn \tx{b}' }		
\end{pchstack}

\vspace{0.1cm}
We say that $\UE$ scheme is $\EUS$-secure if for every PPT-adversary $\adv$ the following advantage is a negligible function of $n$:
\begin{align*}
	\Adv_{\adv}^{\EUS}(n):=\abs{\Pr\left[ \Exp{\UE,\adv}{\EUmod{0}} =1 \right] - \Pr\left[ \Exp{\UE,\adv}{\EUmod{1}} =1 \right]}.
\end{align*}

\vspace{0.3cm}
\noindent
The construction of our $\UE$ scheme is based on $\tx{FrodoPKE}$ learning with errors key encapsulation \cite{Frodo}.

\section{Preliminaries}
For a positive integer $k$, let $[k]:=\{1,\ldots k\}$, and $[k]_{0}:=\{0,1,\ldots k\}$. Denote $\ell_{2}$ and $\ell_{\infty}$ norm by $\|\cdot\|$ and $\|\cdot\|_{\infty}$, respectively. 

Vectors are in column form and they are denoted by bold lower case letters (e.g., $\vr{x}$). We view a matrix as the set of its column vectors and denote by bold capital letters (e.g., $\vr{A}$). The $i$th entry of a vector $\vr{x}$ is denoted $x_{i}$, and the $j$th column of a matrix $\vr{A}$ is denoted $\vr{a}_{j}$ or $\vr{A}[j]$. We identify a matrix $\vr{A}$ with the ordered set $\{\vr{a}_{j}\}$ of its column vector. For convenience, we define the norm of a matrix $\vr{A}$ as the maximum absolute value of its entries, i.e. $\|\vr{A}\|_{\max} = \max_{j}\|\vr{a}_{j}\|_{\infty}$. Note that usually $\|\vr{A}\|_{\max} \neq \|\vr{A} \|_{\infty} $ ($=\sup_{\|\vr{x}\|_{\infty}=1}\|\vr{A}\vr{x}\|_{\infty}$). Roughly speaking, $\|\cdot \|_{\max}$ treats $\vr{A}$ more as a (multi) vector rather than an operator. If the columns of $\vr{A} = \{ \vr{a}_{1}, \ldots ,  \vr{a}_{k} \}$ are linearly independent, then $\vr{A}^{*} = \{ \vr{a}_{1}^{*}, \ldots ,  \vr{a}_{k}^{*} \} $ denote the Gram-Schmidt orthogonalization of vectors $\vr{a}_{1}, \ldots ,  \vr{a}_{k}$ taken in that order. For $\vr{A}\in\RR^{n\times m_{1}}$ and $\vr{B}\in\RR^{n\times m_{2}}$, having an equal number of rows, $[\vr{A}|\vr{B}]\in\RR^{n\times (m_{1}+m_{2})}$ denotes the concatenation of the columns of $\vr{A}$ followed by the columns of $\vr{B}$. Likewise, for  $\vr{A}\in\RR^{n_{1}\times m}$ and $\vr{B}\in\RR^{n_{2}\times m'}$, having an equal number of columns, $[\vr{A};\vr{B}]\in\RR^{(n_{1}+n_{2})\times m}$ is the concatenation of the rows of $\vr{A}$ and the rows of $\vr{B}$.

Let $I$ be a countable set, and let $\Set{X}=\{X_{n}\}_{n\in I}$, $\Set{Y}=\{Y_{n}\}_{n\in I}$ be two families of random variables such that $X_{n}, Y_{n}$ take values in a finite set $\mathcal{R}_{n}$. We call $\Set{X}$ and $\Set{Y}$ \textit{statistically/perfectly indistinguishable} if  their statistical distance is negligible, where the statistical distance between $\{X_{n}\}_{n\in I}$ and $\{Y_{n}\}_{n\in I}$ is defined as the function
%\begin{align*}
	$\Delta(X_{n}, Y_{n}) = \frac{1}{2}\sum_{r\in \mathcal{R}_{n}} \left| \Pr[X_{n}=r] - \Pr[Y_{n}=r]  \right|.$
%\end{align*}
(See \cite{goldreich2001foundations} for more details).

\begin{lemma}[smudging lemma {\normalfont(\cite{asharov2012multiparty})}]\label{lem::smudging}
	Let $B_{1}=B_{1}(n)$, and $B_{2}=B_{2}(n)$ be positive integers and let $e_{1}\in [-B_{1}, B_{1}]$ be a fixed integer. Let $e_{2}\sample [-B_{2}, B_{2}]$ be chosen uniformly at random. If $B_{1}/B_{2}=\negl$, the distribution of $e_{2}$ is statistically indistinguishable from that of $e_{2}+e_{1}$.
\end{lemma}	


\subsection{Lattices}\label{Lattices}
In this subsection we define lattices and present some of their basic properties. An excellent and comprehensive introduction to this theory can be found in \cite{peikert2015decade}.   

\vspace{0.5cm}
Any discrete additive subgroup $\Lat$ of $\RR^{m}$ is called an $m$-dimensional \textit{lattice}, i.e. $\Lat$ is a lattice iff the following conditions hold:
\begin{enumerate}
	\item $\vr{0}\in \Lat$, and if $\vr{v}, \vr{w}\in\Lat$ then $-\vr{v}, \vr{v}+\vr{w}\in \Lat$; 
	\item there is $\varepsilon >0$ such that for every $\vr{v}\in \Lat$ we have $\Lat\cap \{ \vr{w}\mid \| \vr{v}-\vr{w}\|<\varepsilon \} = \{\vr{v}\}$.
\end{enumerate}

Because a lattice $\Lat$ is is an additive subgroup of $\RR^{m}$, it induces the quotient group $\quotient{\RR^{m}}{\Lat}$ of cosets
\begin{align*}
	\vr{x} + \Lat = \{ \vr{x}+\vr{v} \mid \vr{v}\in\Lat  \}, ~~~\vr{x}\in\RR^{m},
\end{align*}
with respect to the addition operation $(\vr{x}+\Lat) + (\vr{y}+\Lat) = (\vr{x}+\vr{y}) + \Lat$.

It turns out that lattices have clearer representation, namely for a lattice $\Lat$ there is a set $\vr{B}$  consisting of linearly independent vectors $\vr{B} = \{\vr{b}_{1}, \ldots , \vr{b}_{k}\}$ such that any lattice point is a integer linear combinations of vectors form $\vr{B}$. 
\begin{align*}
	\Lat = \Lat (\vr{B})=\vr{B}\cdot \ZZ^{k} = \left\{  \sum_{i=1}^{k}\alpha_{i}\vr{b}_{i}\mid \alpha_{i}\in\ZZ  \right\}.
\end{align*}
The set $\vr{B}$ is called a \textit{basis} of $\Lat$, whereas the cardinality $k=\#\vr{B}$ is called the \textit{rank} of $\Lat$.   
If $k=m$ then we say that $\Lat$ is a \textit{full-rank} lattice.  A lattice basis $\vr{B}$ is not unique, namely for any unimodular matrix $\vr{U}\in \ZZ^{m\times m}$ (i.e. $|\det{\vr{U}}|=1$), $\vr{B}\cdot \vr{U}$ is also a basis of $\Lat(\vr{B})$. 

A fundamental domain of $\Lat$ is a connected set $\Fun\subset \RR^{m}$ such that $\vr{0}\in\Fun$ and it contains exactly one representative $\bar{\vr{x}}$ of every coset $\vr{x}+\Lat$. For a lattice $\Lat$ having basis $\vr{B}$, a commonly used fundamental domain is the origin-centered fundamental parallelepiped $\Par(\vr{B}) = \vr{B}\cdot \left( -\frac{1}{2}, \frac{1}{2}  \right]^{m}$, where a coset $\vr{x}+\Lat$ has representative $\vr{x} - \vr{B}\cdot \lfloor \vr{B}^{-1}\cdot \vr{x}\rceil$. The measure of fundamental parallelepiped  can be easily computed as $\vol \Par(\vr{B}) = \sqrt{\det\vr{B}\cdot \vr{B}^{T}}$, in addition this number does not depend on the choice of bases of the lattice, i.e. if $\Lat = \Lat(\vr{B}_{1}) = \Lat(\vr{B}_{2})$ then $\vol \Par(\vr{B}_{1}) = \vol \Par(\vr{B}_{2})$. Therefore, the measure of fundamental parallelepipeds is invariant of the lattice and is called the \textit{determinant} of the lattice $\Lat$ and denoted by $\det\Lat$.

Let $B_{m}(\vr{0}, r):=\{ \vr{x}\in\RR^{m}\mid \|\vr{x}\|< r \}$ be $m$-dimensional open ball of radius $r$ centered at $\vr{0}$. For any $m$-dimensional lattice $\Lat$, the $i$-th minimum $\lambda_{i}(\Lat)$ is the shortest radius $r>0$ such that $B_{m}(\vr{0},r)$ contains $i$ linearly independent lattice vectors
\begin{align*}
	\lambda_{i}(\Lat):=\inf \left\{  r>0\mid \dim(\text{span}(\Lat \cap B_{m}(\vr{0}, r))) \geq i   \right\}.
\end{align*}
For any rank $k$ lattice $\Lat$, $\lambda_{1}(\Lat)\leq \ldots \leq \lambda_{k}(\Lat)$ are constants, and $\lambda_{1}(\Lat) $ is the length of the shortest vector in $\Lat$, e.i. $\lambda_{1}(\Lat) = \min_{\vr{v}\in \Lat/\{\vr{0}\}}\| \vr{v} \|$.

There are several natural computational problems relating to lattices, and for a thorough discussion of the complexity of different types of lattice problems, see \cite{micciancio2002complexity,peikert2015decade}, for example. Due to the overall concept underpinning the paper, we recall the notion of \textit{approximate shortest independent vector problem }($\sivp$).
\begin{definition}\normalfont
	Given a basis $\vr{B}$ of an $m$-dimensional lattice $\Lat = \Lat (\vr{B})$ of the rank $k$, the goal of the \textit{approximate shortest independent vector problem }($\sivp_{\gamma}$) is to find a set of $k$ linearly independent lattice vectors $\vr{V}=\{ \vr{v}_{1}, \ldots , \vr{v}_{k} \}\subset \Lat$, such that $\| \vr{V} \| \leq \gamma{(m)} \cdot \lambda_{k}(\Lat)$, where the approximation factor $\gamma = \gamma(m)$ is a function of the dimension $m$.
\end{definition}
We point out that the most  famous and commonly used version of $\sivp_{\gamma}$ problem concerns the case of full-rank lattices, e.i. when $k=m$. 

\vspace{0.5cm}
A full rank lattice $\Lat$ is called an integer lattice, if $\Lat \subset \ZZ^{m}$, an integer lattice is called a $q$-ary lattice, if $q\ZZ^{m}\subset \Lat \subset \ZZ^{m}$, where $q\in\ZZ_{\geq 1}$. 

By definition, a lattice $\Lat=\Lat(\vr{B})$ is an integer lattice iff $\vr{B}\in \ZZ^{m\times m}$ is an integer square matrix, so the determinant $\det \Lat$ is a positive integer. It easy to show that an integer lattice is always a ($\det \Lat$)-ary lattice. For a $q$-ary lattice $\Lat$ the following simple conclusions in algebra holds:
\begin{itemize}
	\item $\quotient{\ZZ^{m}}{q\ZZ^{n}} \cong \ZZ_{q}^{m}$ (additive group isomorphism);
	\item $\quotient{\ZZ^{m}}{\Lat} \cong \quotient{\ZZ_{q}^{m}}{(\quotient{\Lat}{q\ZZ^{m}})}$ (additive group isomorphism).
\end{itemize}
In addition, $\quotient{\ZZ^{m}}{\Lat} $ is a finite group and $|\quotient{\ZZ^{m}}{\Lat} | = \det \Lat$.

\subsection{Gaussians}\label{Gaussian}
For any real $s>0$ and $\vr{c}\in\RR^{m}$, the Gauss function $\rho_{s,\vr{c}}$ centered on $\vr{c}$ with parameter $s$ is defined as
\begin{align*}
	\rho_{s,\vr{c}}(\vr{x})  = \exp \left(-\frac{\pi}{s^{2}}\|\vr{x} - \vr{c} \|^{2}\right), ~~~\vr{x}\in\RR^{m},
\end{align*}
and 
\begin{align*}
	\rho_{s} = \rho_{s, \vr{0}}, ~~~\rho = \rho_{1},~\text{i.e.}~ \rho(\vr{x}) = e^{-\pi\| \vr{x} \|^{2}}.  
\end{align*}
Directly from the definition we have $\rho_{s}(\vr{x}) = \rho\left(  {s}^{-1} \vr{x}\right)$.

For any real $s>0$, \textit{the rounded Gaussian distribution} with parameter (or width) $s$, denoted $\Psi_{s}$, is the distribution over $\ZZ$ obtained by rounding a sample from $\Distr_{s}$ to the nearest integer:
\begin{align*}
	\Psi_{s}(x)= \int_{\{ z\mid \lfloor z \rceil = x\}} \Distr_{s}(z)dz.
\end{align*}

Let $\Lat\subseteq \ZZ^{m}$ be a lattice and let $\rho_{s, \vr{c}} (\Lat) = \sum_{\vr{x}\in \Lat} \rho_{s, \vr{c}}(\vr{x})$. Define the discrete Gaussian distribution over $\Lat$ with center $\vr{c}$, and parameter $s$ as
\begin{align*}
	\Distr_{\Lat, s, \vr{c}} (\vr{x}) =  \frac{\rho_{s, \vr{c}}(\vr{x})}{\rho_{s, \vr{c}}(\Lat)}, ~~~\vr{x}\in\Lat.
\end{align*}
For notational convenience, we let
\begin{align*}
	\Distr_{\Lat, s} = \Distr_{\Lat, s, \vr{0}}, ~ \Distr_{s, \vr{c}}^{m} = \Distr_{\ZZ^{m}, s, \vr{c}}, ~ \Distr_{s}^{m} = \Distr_{\ZZ^{m}, s}.
\end{align*}

%\begin{lemma}{\normalfont (\cite{banaszczyk1993new,lindner2011better})}
%	For any $s,t\in\RR_{>0}$, $c\in\RR_{\geq1}$, $C=c\cdot \exp{\left( 2^{-1}\cdot(1-c^{2}) \right)}$, $m\in\ZZ_{>0}$, and any $\vr{y}\in\RR^{m}$ the following conditions hold: 
%	\begin{itemize}
	%		\item $\Pr_{\vr{x}\gets  \Distr_{s}^{m}} \left[ \| \vr{x}\|_{\infty}> t\cdot s \right]\leq 2e^{-\pi t^{2}}$;
	%		\item $\Pr_{\vr{x}\gets  \Distr_{s}^{m}} \left[ \| \vr{x}\|_{\infty}>   c\cdot \frac{1}{\sqrt{2\pi}} \cdot s\sqrt{m}   \right]\leq C^{m}$;
	%		\item $\Pr_{\vr{x}\gets  \Distr_{s}^{m}} \left[    \vert  \inner{x}{y} \vert  > t\cdot s\norm{y} \right]         \leq 2e^{-\pi t^{2}}$.
	%	\end{itemize}
%\end{lemma}

\vspace{0.5cm}
We summarize several facts from the literature about discrete Gaussians over lattices, again specialized to our family of interest.
\begin{lemma}\label{lem::02}Let $n<m$ and $\vr{T}$ be any basis of $\Lat$ for some $\vr{A}\in\ZZ_{q}^{n\times m}$ whose columns generate $\ZZ_{q}^{n}$, let $\vr{u}\in\ZZ_{q}^{n}$ and $\vr{c}\in\ZZ^{m}$
	be arbitrary, and let $s\geq \| \vr{T}^{*} \| \cdot \omega \left( \sqrt{\log{m}} \right)$, we have
	\begin{enumerate}
		\item {\normalfont \cite{micciancio2007worst,cash2012bonsai}:}
		$\Pr_{\vr{x}\gets  \Distr_{\Lat}} \left[ \| \vr{x}-\vr{c}\| >  s \cdot \sqrt{m} \right]\leq \negl.$
		
		\item {\normalfont \cite{peikert2006efficient,cash2012bonsai}:} 
		$\Pr_{\vr{x}\gets  \Distr_{\Lat}} \left[ \vr{x}=0 \right]\leq \negl.$
		
	\end{enumerate}
\end{lemma}


\section{Frodo based UE scheme}
\subsection{Frodo KEM}\label{sec::Frodo}
Since the construction of an $\UE$ scheme is based on $\tx{FrodoPKE}$, below we present a framework for a learning with errors key encapsulation $\tx{Frodo}$ \cite{Frodo,Frodo2}. To this end, let $D\in\ZZ_{>0}$ and $q=2^{D}$, $B\in\ZZ_{[0,D]}$, $n,\bar{m}, \bar{n}\in \ZZ_{>0}$, where $n\equiv 0 \pmod{8}$. In addition, $\Set{M}=\{0,1\}^{\ell}$ (see below) denote the messages space. 

The error distribution $\chi$ used in $\tx{FrodoPKE}$ is a discrete, symmetric distribution on $\ZZ$, centered at zero and with small support, which approximates a rounded continuous Gaussian distribution. The support of $\chi$ in $S_{\chi}=\{ -s, -s+1, \ldots -1, 0, 1, \ldots , s-1, s \}$ for $s\in\ZZ_{>0}$. The probabilities $\chi(z) = \chi(-z)$ for $z\in S_{\chi}$ are given by a discrete probability density function, which is described by a table $T_{\chi} = (T_{\chi}(0), T_{\chi}(1), \ldots , T_{\chi}(s))$ of $s+1$ positive integers related to the cumulative distribution function. Sampling from $\chi$ is done via inversion sampling methods, see \cite{Frodo,Frodo2} for details.
The distribution $\chi$ is used as an integral part of the algorithm $\tx{Frodo.SampleMatrix}$ (see \cite{Frodo,Frodo2}) to sample error matrices. However, to ease notation, the outputs from this algorithm will be denoted by $\vr{E}\sample \chi(\ZZ_{q}^{\dim})$. 

$\tx{FrodoPKE}$ is defined as a tuple of PPT-algorithms $(\tx{FrodoPKE.Setup},$ $\tx{FrodoPKE.KeyGen},$\\ $\tx{FrodoPKE.Enc},$ $\tx{FrodoPKE.Dec})$  that have the following form:

\vspace{0.5cm}
 \begin{pchstack}[center]
	
	\begin{pcvstack}
		
		\procedure{$\tx{FrodoPKE.Setup}(1^{n})$}{
			\pcln \vr{A}\sample \ZZ_{q}^{n\times n}\\
			\pcln \params\gets (\vr{A}, 1^{n},D, B, n,\bar{m}, \bar{n})\\
			\pcln \pcreturn \params}
		
		\pcvspace
		
		\procedure{$\tx{FrodoPKE.KeyGen}(\params)$}{
			\pcln \vr{S}, \vr{E} \sample \chi(\ZZ_{q}^{n\times \bar{n}})\\
			\pcln \vr{B} \gets \vr{A}\cdot \vr{S} + \vr{E} \in \ZZ_{q}^{n\times \bar{n}}\\
			\pcln \pcreturn \pk\gets \vr{B}, ~\sk\gets \vr{S}}
		
		
	\end{pcvstack}
	
	\pchspace
	
	\begin{pcvstack}
		~~~~~
     \end{pcvstack}	
		
	\pchspace	
	
	\begin{pcvstack}
	
	\procedure{$\tx{FrodoPKE.Enc}(\pk, \m\in \Set{M})$}{
		\pcln \vr{S}', \vr{E}' \sample \chi(\ZZ_{q}^{\bar{m}\times {n}})\\
		\pcln \vr{E}'' \sample \chi(\ZZ_{q}^{\bar{m}\times \bar{n}})\\
		\pcln \vr{B}' \gets \vr{S}' \cdot\vr{A} + \vr{E}' \in \ZZ_{q}^{\bar{m}\times {n}}\\
		\pcln \vr{V} \gets \vr{S}' \cdot\vr{B} + \vr{E}'' \in \ZZ_{q}^{\bar{m}\times \bar{n}}\\
		\pcln  \vr{C}_{1} \gets \vr{B}', ~\vr{C}_{2} \gets \vr{V} + \tx{encode}(\m)\\
		\pcln \pcreturn \ct\gets(\vr{C}_{1}, \vr{C}_{2})}
	
	\pcvspace  
	
	\procedure{$\tx{FrodoPKE.Dec}(\ct, \sk)$}{
		\pcln \vr{M} \gets \vr{C}_{2} - \vr{C}_{1}\cdot \vr{S}  \in \ZZ_{q}^{\bar{m}\times \bar{n}}\\
		\pcln \m'\gets \tx{decode}(\vr{M})\\
		\pcln\pcreturn \m'
	}
\end{pcvstack}	
\end{pchstack}	

 \vspace{0.5cm}
Although the functions $\tx{encode}$ and $\tx{decode}$ are formally defined in \cite{Frodo,Frodo2}, we briefly describe the idea standing behind their constructions. The function $\tx{encode}$ encodes bit strings of length $\ell = B\cdot\bar{m}\cdot\bar{n}$ as $\bar{m}\times\bar{n}$ matrices over $\ZZ_{q}$. To this end, a subsequent $B$-bit sub-string is treated as an integer $k\in[0,2^{B}]$ and transformed by $\tx{ec}(k)=k\cdot q/2^{B}$ into element of $\ZZ_{q}$. The corresponding $\tx{decode}$ function reverses this operation by taking advantage of $\tx{dc}(c)=\lfloor c\cdot 2^{B}/q\rceil \pmod 2^{B}$. 
It turns out that the scheme is correct if an error is sufficiently small.
\begin{lemma}(\cite{Frodo})\label{lem::01}
If $q=2^{D}$ and $B\leq D$, then $\tx{dc}(\tx{ec}(k)+e)=k$ for any integer $0\leq k < 2^{B}$ and $-q/2^{B+1}\leq e < q/2^{B+1}$. 	
\end{lemma}	


\section{Construction of  EU scheme based on Frodo}
%In the introduction we explained that it is natural to see a token $\tok_{\epoch+1}$ as a homomorphic encryption of a secret key $\sk_{\epoch}$ under a public key $\pk_{\epoch+1}$. 
Note that to update a ciphertext $\ct_{\epoch}$, we must homomorphically decrypt $\ct_{\epoch}$ by using a token $\tok_{\epoch+1}$ and obtain $\ct_{\epoch+1}\gets \tx{FrodoPKE.Enc}(\pk_{\epoch+1}, \m)$.
We recall that it is natural to see a token $\tok_{\epoch+1}$ as a homomorphic encryption of a secret key $\sk_{\epoch}$ under a public key $\pk_{\epoch+1}$. Functions \textit{bit-ordering} ($\tx{Ord}$) and \textit{Tensor-$D$} ($\bigotimes_{D}$) that make it possible to meet these requirements are formally defined below. The idea behind this approach is drawn from the key-switching technique \cite{brakerski2014efficient,brakerski2014leveled,brakerski2011fully} and the noise smudging lemma \cite{asharov2012multiparty, nishimaki2022direction}.

\begin{description}
	
	\item[{Bit-ordering}] {($\tx{Ord}$)}. This function takes as input a matrix $\vr{X}\in \ZZ_{q}^{\bar{m}\times \bar{n}}$ and outputs $\vr{Y}\in \ZZ_{q}^{\bar{m}\times \bar{n}D}$. If $\vr{X}[i,\cdot]=[{x}_{i1}, {x}_{i2}, \ldots , {x}_{i\bar{n}}]$ is the $i$-th row of $\vr{X}$, then each entry ${x}_{i,j}\in\ZZ_{q}$ is encoded on $D$-bits, i.e. ${x}_{i,j}=\sum_{k=1}^{D}2^{k-1}y_{i,j,k}$. $\tx{Ord}$ transforms $\vr{X}[i,\cdot]$ into $i$-th row $\vr{Y}[i,\cdot]=[\vr{y}_{i1}, \vr{y}_{i2}, \ldots , \vr{y}_{iD}]$ of $\vr{Y}$, where $\vr{y}_{i,k}=[y_{i,j,k}]_{j}\in\bin^{\bar{n}}$.
	
	\item[Tensor-$D$] ($\bigotimes_{D}$  ). This function transforms $\ZZ_{q}^{n\times \bar{n}}$ into $\ZZ_{q}^{nD\times \bar{n}}$. Let $\vr{X} =[\vr{x}_{1}, \vr{x}_{2}, \ldots , \vr{x}_{\bar{n}}]\in \ZZ_{q}^{n\times \bar{n}}$, where $\vr{x}_{j}$ is the $j$-th column of $\vr{X}$. Then $\bigotimes_{D}(\vr{X}) = \vr{Y}=[\vr{y}_{1}, \vr{y}_{2}, \ldots , \vr{y}_{\bar{n}}]\in \ZZ_{q}^{nD\times \bar{n}}$, where the $j$-th column $\vr{y}_{j}$ is of the form $\vr{y}_{j}=[1, 2, \ldots , 2^{D-1}]^{T}\otimes \vr{x}_{j}$.
	
	\item[Formal product] Let $\vr{A}\in\ZZ_{q}^{l\times r}$, $\vr{B}\in\ZZ_{q}^{r\times s_{1}}$, and $\vr{C}\in\ZZ_{q}^{r\times s_{2}}$. The formal product $\vr{A}\cdot (\vr{B},\vr{C})$ should be viewed as a product $\vr{A}\cdot [\vr{B}\mid \vr{C}]$, that equals $[\vr{A}\cdot \vr{B}\mid \vr{A}\cdot \vr{C}]$. Now, splitting the last matrix 'along' the concatenation operator, we get the formal condition $(\vr{A}\cdot \vr{B},  \vr{A}\cdot \vr{C}) = \vr{A}\cdot (\vr{B},\vr{C})$. 
\end{description}

\begin{remark}
Note that the following equality holds
$\tx{Ord}(\vr{C}_{\epoch , 1})\cdot \bigotimes_{D} (\vr{S}_{\epoch}) = \vr{C}_{\epoch , 1}\cdot \vr{S}_{\epoch}$.
\end{remark}	


Having defined the functions \textit{bit-ordering}, \textit{tensor-$D$} and \textit{formal product}, we are ready to provide a description of our UE scheme:

\vspace{0.5cm}
 \begin{pchstack}[center]	
	\begin{pcvstack}
		\procedure{$\tx{UE.Setup}(1^{n})$}{
			\pcln \params\sample \tx{FrodoPKE.Setup}(1^{n})\\
			\pcln \params=(\vr{A}, 1^{n},D, B, n,\bar{m}, \bar{n})\\
			\pcln\pcreturn \params}
		
		\pcvspace
		
		\procedure{$\uekg(\params)$}{
			\pcln (\vr{B}_{\epoch}, \vr{S}_{\epoch})  \sample \tx{FrodoPKE.KeyGen}(\params)\\
			\pcln \sk_{\epoch}\gets \vr{S}_{\epoch},~\pk_{\epoch} \gets \vr{B}_{\epoch}\\
			\pcln\pcreturn \key_{\epoch} \gets (\sk_{\epoch}, \pk_{\epoch})}
		
	\end{pcvstack}
	
	\pchspace
	\hspace{0.7cm}
	
	\begin{pcvstack}
		\procedure{$\ueenc(\key_{\epoch}, \m\in\Set{M})$}{
			\pcln   \key_{\epoch} = (\vr{S}_{\epoch}, \vr{B}_{\epoch}) \\
			\pcln (\vr{C}_{1}, \vr{C}_{2}) \sample \tx{FrodoPKE.Enc}(\vr{B}_{\epoch}, \m)\\
			\pcln  \ct\gets(\vr{C}_{1}, \vr{C}_{2}) \in  \ZZ_{q}^{\bar{m}\times {n}}  \times \ZZ_{q}^{\bar{m}\times \bar{n}}\\
			\pcln\pcreturn  \ct}
		
		\pcvspace
		
		\procedure{$\uedec(\key_{\epoch}, \ct)$}{
			\pcln   \key_{\epoch} = (\vr{S}_{\epoch}, \vr{B}_{\epoch}), ~\ct = (\vr{C}_{1}, \vr{C}_{2})  \\
			\pcln \m' \gets\tx{FrodoPKE.Dec}(\ct, \vr{S}_{\epoch})\\
			\pcln\pcreturn  \m'}
	\end{pcvstack}
\end{pchstack}	

\begin{pcvstack}[center]	
	\procedure{$\uetg(\key_{\epoch}, \key_{\epoch+1})$}{
		\pcln   \key_{\epoch} = (\vr{S}_{\epoch}, \vr{B}_{\epoch}), ~  \key_{\epoch+1} = (\vr{S}_{\epoch+1}, \vr{B}_{\epoch+1}) \\
		\pcln  \vr{S}_{(1)}', \vr{E}_{(1)}' \sample \chi(\ZZ_{q}^{nD\times {n}}), ~ \vr{E}_{(1)}'' \sample \chi(\ZZ_{q}^{nD\times \bar{n}}) \\
		\pcln \vr{S}_{(2)}', \vr{E}_{(2)}' \sample \chi(\ZZ_{q}^{n\times {n}}), ~ \vr{E}_{(2)}'' \sample \chi(\ZZ_{q}^{n\times \bar{n}}) \\
		\pcln \tok_{\epoch+1}^{(1)} \gets \left( \vr{S}_{(1)}'\cdot\vr{A}  +  \vr{E}_{(1)}',  \vr{S}_{(1)}'\cdot\vr{B}_{\epoch+1}+\vr{E}_{(1)}''  - {{\scriptsize\bigotimes}_{D}}(\vr{S}_{\epoch}) \right) \\
		\pcln \tok_{\epoch+1}^{(2)} \gets \left( \vr{S}_{(2)}'\cdot\vr{A}  +  \vr{E}_{(2)}',  \vr{S}_{(2)}'\cdot\vr{B}_{\epoch+1}+\vr{E}_{(2)}''\right)    \\
		\pcln\pcreturn  \tok_{\epoch+1}\gets (  \tok_{\epoch+1}^{(1)} ,  \tok_{\epoch+1}^{(2)} )}
\end{pcvstack}		

\vspace{0.5cm}
%%%%

\begin{pcvstack}[center]	
	%\begin{pcvstack}[center]		
	\procedure{$\ueupd( \tok_{\epoch+1}, \ct_{\epoch})$}{
		\pcln  \tok_{\epoch+1} = (  \tok_{\epoch+1}^{(1)} ,  \tok_{\epoch+1}^{(2)} ), ~ \ct_{e}=(\vr{C}_{\epoch,1}, \vr{C}_{\epoch,2}) ~~~~~~~~~~~~~~~~~~~~~~~~~~~~~~\\
		\pcln (\vr{C}_{1}^{(1)},  \vr{C}_{2}^{(1)}) \gets \tx{Ord}(\vr{C}_{\epoch,1})\cdot \tok_{\epoch+1}^{(1)}\\
		\pcln \vr{R} \sample \chi(\ZZ_{q}^{\bar{m}\times {n}})\\
		\pcln (\vr{C}_{1}^{(2)},  \vr{C}_{2}^{(2)}) \gets \vr{R}\cdot  \tok_{\epoch+1}^{(2)}\\
		\pcln \ct_{\epoch+1}\gets \left(  \vr{C}_{1}^{(1)} + \vr{C}_{1}^{(2)}, \vr{C}_{\epoch,2} + \vr{C}_{2}^{(1)} + \vr{C}_{2}^{(2)}\right)\\
		\pcln\pcreturn  \ct_{\epoch+1} }
	
\end{pcvstack}		

\vspace{0.5cm}
\begin{remark}\label{rem::01}
	Note that in order to generate a token $\tok_{\epoch+1}$, $\uetg$ requires merely $\sk_{\epoch}$ and $\pk_{\epoch+1}$ instead of ,,full'' keys $\key_{\epoch}, \key_{\epoch+1}$. This is consistent with the idea discussed in the introduction.  
\end{remark}	

The conditions that make the scheme correct are provided in the next theorem. 
%The scheme is correct. Indeed, 
\begin{theorem}
	Let $\chi$ be an error distribution defined in Sec. \ref{sec::Frodo} such that $S_{\chi}=\{ -s, \ldots, 0,  \ldots ,s \}$ is its support, and
	let $T$ be the maximum number of the epochs. If $2(n^{2}Ds^{3}+n^{2}s^{3})+nDs + ns^{2}<q/(T\cdot 2^{B+1})$, then the scheme is correct.
\end{theorem}	


\begin{proof}
Let $\ct_{\epoch} =  (\vr{C}_{\epoch,1}, \vr{C}_{\epoch, 2})$ and $\ct_{\epoch+1} =  (\vr{C}_{\epoch+1,1}, \vr{C}_{\epoch+1, 2})$, where the coefficients are as in $\ueupd$. It is easy to check that 
 \begin{align*}
	\vr{C}_{\epoch+1,2} - \vr{C}_{\epoch +1, 1}\cdot \vr{S}_{\epoch +1} = \vr{C}_{\epoch,2} - \vr{C}_{\epoch , 1}\cdot \vr{S}_{\epoch } + \tx{Error}
\end{align*}
where
\begin{align}\label{pr01::01}
	\tx{Error}= \left(\tx{Ord}(\vr{C}_{\epoch , 1})\cdot\vr{S}_{(1)}' + \vr{R}\cdot \vr{S}_{(2)}'\right)\vr{E}- \left(\tx{Ord}(\vr{C}_{\epoch , 1})\cdot\vr{E}_{(1)}' + \vr{R}\cdot \vr{E}_{(2)}'\right)\vr{S}_{\epoch+1} \\
	+ \tx{Ord}(\vr{C}_{\epoch , 1})\cdot\vr{E}_{(1)}'' + \vr{R}\cdot \vr{E}_{(2)}'' \nonumber
\end{align}
Since $\tx{Ord}(\vr{C}_{\epoch , 1})\in\{0,1\}^{\bar{m}\times nD}$, $\vr{S}_{(1)}' \in S_{\chi}^{nD\times n}$, and $\vr{E} \in S_{\chi}^{n\times n}$, we have
\begin{align}\label{pr01::02}
	\left\|  \left(\tx{Ord}(\vr{C}_{\epoch , 1})\cdot\vr{S}_{(1)}' + \vr{R}\cdot \vr{S}_{(2)}'\right)\vr{E} \right\|_{\max}\leq n^{2}Ds^{2}+n^{2}s^{3}.
\end{align}
Similarly, we obtain the following bounds
\begin{align}
	\left\|  \left(\tx{Ord}(\vr{C}_{\epoch , 1})\cdot\vr{E}_{(1)}' + \vr{R}\cdot \vr{E}_{(2)}'\right)\vr{S}_{\epoch+1}    \right\|_{\max}\leq n^{2}Ds^{2}+n^{2}s^{3}, \label{pr01::03}\\
		\left\|  \tx{Ord}(\vr{C}_{\epoch , 1})\cdot\vr{E}_{(1)}'' \right\|_{\max}\leq nDs, ~~~ \left\|  \vr{R}\cdot \vr{E}_{(2)}'' \right\|_{\max}\leq ns^{2}. \label{pr01::04}
\end{align}
By (\ref{pr01::01}), triangle inequality and (\ref{pr01::02}) -(\ref{pr01::04}), we have
\begin{align*}
	\norm{\tx{Error}}_{\infty}\leq 2(n^{2}Ds^{3}+n^{2}s^{3})+nDs + ns^{2}<q/(T\cdot 2^{B+1}).
\end{align*}
Therefore, when we update a ciphertext once, an noise is added and its magnitude is at most $q/(T\cdot 2^{B+1})$. This means that if we update a ciphertext $T$ times, the total noise does not exceed $q/2^{B+1}$. Hence, by Lemma \ref{lem::01} the proof is complete.
\end{proof}	

\subsection{Security proof}
In this section we justify that the scheme is $\EUS$-secure in the backward-leak uni-directional setting. To this end, we follow the firewall technique \cite{jiang2020,nishimaki2022direction}.

Let us start with the explanation that the scheme meets, in particular, uni-directional ciphertext updates from its security.
\begin{theorem}
	The advantage of any PPT adversary $\adv$ in converting a ciphertext under a public key $\pk_{\epoch+1}$ into a ciphertext under $\pk_{\epoch}$ by using $(\sk_{\epoch},\pk_{\epoch}), \pk_{\epoch+1}$ and $\tok_{\epoch+1}$ is negligible.
\end{theorem}	

\begin{proof}
Suppose, for contradiction, that $\adv$ can convert a ciphertext under $\pk_{\epoch+1}$ into one under $\pk_{\epoch}$. Since $\tx{FrodoPKE}$ is $\indcpa$-secure, we construct a PPT adversary $\adversary{B}$ that is able to break the security under $\pk_{\epoch+1}$.  

Algorithm $\adversary{B}$ is given $\params$ and $\pk_{\epoch+1}$. $\adversary{B}$ generates $(\sk_{\epoch}, \pk_{\epoch})\gets \uekg(\params)$ and $\tok_{\epoch+1}\gets\uetg(\sk_{\epoch},\pk_{\epoch+1})$ (see Remark \ref{rem::01}), 
computes $\m_{0}, \m_{1}\in\Set{M}$, of the same length, and sends them to its challenger. The challenger computes $\ct^{*}\gets \tx{FrodoPKE.Enc}(\pk, \m_{b})$, where $b\sample\bin$, and sends $\ct^{*}$ to $\adversary{B}$. Next $\adversary{B}$ run $\adv$ on input $((\sk_{\epoch},\pk_{\epoch}),\tok_{\epoch+1}, \ct^{*})$. $\adv$ outputs a ciphertext $\ct'$ under $\pk_{\epoch}$. Eventually, $\adversary{B}$ computes $\m'\gets \tx{FrodoPKE.Dec}(\ct', \sk_{\epoch})$, and if $\m' = \m_{b'}$ it outputs $b'$. Obviously, $b'=b$ with overwhelming probability. This completes the proof.

\end{proof}	

In the next step we reorganize the security game and make new, hybrid, algorithms that simulate the scheme up to negligible probability and at the same time are more convenient to analyze. To this end, let us look again at the update procedure. By definition of the scheme, we have 

\begin{align*}
	 \vr{C}_{1}^{(1)} &= \tx{Ord}(\vr{C}_{\epoch , 1})\cdot\vr{S}_{(1)}'  \cdot \vr{A} + \tx{Ord}(\vr{C}_{\epoch , 1})\cdot \vr{E}_{(1)}' ,     \\
	 %
	 \vr{C}_{2}^{(1)} &=   \tx{Ord}(\vr{C}_{\epoch , 1})\cdot \vr{S}_{(1)}' \cdot \vr{B}_{\epoch+1} + \tx{Ord}(\vr{C}_{\epoch , 1})\cdot \vr{E}_{(1)}''  -   \underbrace{\tx{Ord}(\vr{C}_{\epoch , 1})\cdot \bigotimes_{D} (\vr{S}_{\epoch}) }_{\vr{C}_{\epoch , 1}\cdot \vr{S}_{\epoch}}, \\
	 %
	 	\vr{C}_{1}^{(2)} &= \vr{R}\cdot \vr{S}_{(2)}' \cdot \vr{A} + \vr{R}\cdot \vr{E}_{(2)}'  ,\\
	 	%
	 \vr{C}_{2}^{(2)} &= \vr{R}\cdot \vr{S}_{(2)}' \cdot \vr{B}_{\epoch+1} + \vr{R}\cdot \vr{E}_{(2)}''.
\end{align*}
By simple calculation, we get
\begin{align*}
 \vr{C}_{1}^{(1)} + \vr{C}_{1}^{(2)}  = \left(  \underbrace{\tx{Ord}(\vr{C}_{\epoch , 1})\cdot \vr{S}_{(1)}'  +  \vr{R}\cdot \vr{S}_{(2)}' }_{\vr{S}^{\dagger}}  \right) \cdot \vr{A}  + \underbrace{\tx{Ord}(\vr{C}_{\epoch , 1})\cdot \vr{E}_{(1)}' + \vr{R}\cdot \vr{E}_{(2)}' }_{\vr{E}^{\dagger}} ,
\end{align*}
and
\begin{align*}
 \vr{C}_{2}^{(1)}  +  \vr{C}_{2}^{(2)}  + \vr{C}_{\epoch,2}  = & \left(  \underbrace{\tx{Ord}(\vr{C}_{\epoch , 1})\cdot \vr{S}_{(1)}'  +  \vr{R}\cdot \vr{S}_{(2)}' }_{\vr{S}^{\dagger}}  \right) \cdot \vr{B}_{\epoch+1} \\
 & + \underbrace{\tx{Ord}(\vr{C}_{\epoch , 1})\cdot \vr{E}_{(1)}'' + \vr{R}\cdot \vr{E}_{(2)}'' +\vr{E}_{\epoch}''}_{\vr{E}^{\dagger\dagger}} \\
 &- \underbrace{\vr{E}'\cdot \vr{S}_{\epoch} + \vr{S}' \cdot \vr{E}}_{\text{''garbage''}}~ +  ~\tx{encode}(\m)
\end{align*}
%
This leads us to the following conclusion.
\begin{concl}\label{conc::01}
	By Lemma \ref{lem::smudging},  a ciphertext $\ct_{\epoch+1}$ is statistically indistinguishable from $(\vr{S}^{\dagger}\cdot\vr{A} + \vr{E}^{\dagger},  \vr{S}^{\dagger}\cdot\vr{B}_{\epoch+1} + \vr{E}^{\dagger\dagger} +  \tx{encode}(\m))$. This means that we can simulate an update ciphertext by using the original ciphertext, its plaintext and randomness, the new epoch public key, and randomness for generating the token $\tok_{\epoch+1}$.
\end{concl}

\vspace{0.5cm}
%\begin{pcvstack}[center]	

\procedure{$\tx{Hyb.UE.Upd}(\ct_{\epoch}, \vr{B}_{\epoch+1},\m, \vr{E}_{\epoch}'', (\vr{S}_{(1)}', \vr{S}_{(2)}', \vr{E}_{(1)}'', \vr{E}_{(2)}''))$}{
	\pcln \text{Parse} ~~ \ct_{\epoch} = (\vr{C}_{\epoch, 1}, \vr{C}_{\epoch,2})\\
	\pcln \text{Choose} ~~ \vr{R} \sample \chi(\ZZ_{q}^{\bar{m}\times {n}})\\
	\pcln \text{Set} ~~ \vr{S}^{\dagger} = \tx{Ord}(\vr{C}_{\epoch , 1})\cdot \vr{S}_{(1)}'  +  \vr{R}\cdot \vr{S}_{(2)}' \\
	\pcln \text{Set} ~~ \ct_{\epoch+1} =  \left(\vr{S}^{\dagger}\cdot\vr{A} + \vr{E}^{\dagger},  \vr{S}^{\dagger}\cdot\vr{B}_{\epoch+1} + \vr{E}^{\dagger\dagger} +  \tx{encode}(\m)\right) \\
	\pcln \text{Output} ~~ \left( \ct_{\epoch+1} ,  \vr{E}_{\epoch}''  \right).
}

\vspace{0.3cm}

\procedure{$\tx{Sim.UE.KG}(\params)$}{
	\pcln \text{Choose} ~~ \vr{B}_{e}^{+} \sample  \ZZ_{q}^{n\times \bar{n}} \\
	\pcln \text{Output} ~~ \pk_{\epoch} = \vr{B}_{e}^{+}.}


\vspace{0.3cm}


\procedure{$\tx{Sim.UE.TG}(\params)$}{
	\pcln \text{Choose} ~~\tok_{\epoch+1}^{(1)+} \sample  \ZZ_{q}^{nD\times {n}} \times \ZZ_{q}^{nD\times \bar{n}},~  \text{and} ~ \tok_{\epoch+1}^{(2)+}  \sample    \ZZ_{q}^{n\times {n}} \times \ZZ_{q}^{D\times \bar{n}}  \\
	\pcln \text{Output} ~~ \tok_{\epoch+1}^{+} = \left( \tok_{\epoch+1}^{(1)+},   \tok_{\epoch+1}^{(2)+}  \right).}

\vspace{0.3cm}


\procedure{$\tx{Sim.UE.Upd}(\params)$}{
	\pcln \text{Choose} ~~  (\vr{C}_{1}, \vr{C}_{2}) \sample  \ZZ_{q}^{\bar{m}\times {n}}  \times \ZZ_{q}^{\bar{m}\times \bar{n}} \\
	\pcln \text{Output} ~~ \ct_{\epoch+1} = (\vr{C}_{1}, \vr{C}_{2}).}


\vspace{0.3cm}

\procedure{$\tx{Sim.UE.Enc}(\params)$}{
	\pcln \text{Choose} ~~  (\vr{C}_{1}, \vr{C}_{2}) \sample  \ZZ_{q}^{\bar{m}\times {n}}  \times \ZZ_{q}^{\bar{m}\times \bar{n}} \\
	\pcln \text{Output} ~~ \ct_{\epoch} = (\vr{C}_{1}, \vr{C}_{2}).}

%\end{pcvstack}	(\vr{C}_{1}, \vr{C}_{2}) \in  \ZZ_{q}^{\bar{m}\times {n}}  \times \ZZ_{q}^{\bar{m}\times \bar{n}}

\vspace{1cm}

\noindent Directly from Conclusion \ref{conc::01}, we infer that the following lemma holds.
\begin{lemma}
	$\ueupd( \tok_{\epoch+1}, \ct_{\epoch})$ is statistically indistinguishable form $\tx{Hyb.UE.Upd}(\ct_{\epoch}, \vr{B}_{\epoch+1},\m,$ $ \vr{E}_{\epoch}'', (\vr{S}_{(1)}', \vr{S}_{(2)}', \vr{E}_{(1)}'', \vr{E}_{(2)}''))$.	
\end{lemma}	

\noindent By the above lemma, we get the following conclusion. 

\begin{concl}\label{conc::02}
	The oracle $\Upd(\ct_{\epoch})$ can be simulated by $\tx{Hyb.UE.Upd}(\ct_{\epoch}, \vr{B}_{\epoch+1},\m, \vr{E}_{\epoch}'', (\vr{S}_{(1)}',$ $\vr{S}_{(2)}', \vr{E}_{(1)}'', \vr{E}_{(2)}''))$.	
\end{concl}



Let $T+1$ be a number of epochs (they are counted from zero). Define the following hybrid games: 
\begin{description}
	\item[$\tx{Hyb}_{i}(b):$] This is the same as $\Exp{\UE,\adv}{\EUmod{b}}(n)$ except the following difference: When the adversary sends a query $(\bar{\m}, \bar{\ct})$ to $\Chall$ or an empty query to $\UpdC$ at epoch $j$:
	\begin{itemize}
		\item for $j<i$ return an honestly generated challenge-equal ciphertext, i.e. 
		\begin{itemize}
			\item[{}] \textbf{if} $b=0$ \textbf{then} $\ueenc(\key_{\widetilde{\epoch}}, \bar{\m})$,
			\item[{}] \textbf{else} $\tx{UE.Upd}(\tok_{\widetilde{\epoch}}, \bar{\ct})$. 
		\end{itemize}
		\item $j\geq i$, return a random ciphertext.
	\end{itemize}
\end{description}
It is easy to see that $\tx{Hyb}_{T+1}(b)$ is the same as the original $\tx{rand\md ind\md eu\md cpa}$	game in the backward-leak uni-directional setting $\Exp{\UE,\adv}{\EUmod{b}}(n)$. Let $\mathcal{U} (n)$ be a random variable distributed uniformly in $[0,T]$, by the standard hybrid argument, we have
\begin{align}
	\Adv_{\adv}&^{\EUS}(n) = \left| \Pr\left[ \tx{Hyb}_{T+1}(1)=1 \right]   -  \Pr\left[ \tx{Hyb}_{T+1}(0)=1 \right]   \right| \nonumber\\
	=&\left|  \sum_{i=0}^{T} \left\{   \left(  \Pr\left[   \tx{Hyb}_{i+1}(1)=1 \right]   -   \Pr\left[   \tx{Hyb}_{i}(1)=1 \right] \right) \right. \right. \nonumber\\
	& -   \left.\left. \left(  \Pr\left[   \tx{Hyb}_{i+1}(0)=1 \right]   -   \Pr\left[   \tx{Hyb}_{i}(0)=1 \right] \right)   \right\} \right|    \label{eq::02}\\
	%
	=&\left|  \sum_{i=0}^{T}  \left(  \Pr\left[   \tx{Hyb}_{\mathcal{U}(n)+1}(1)=1 \mid \mathcal{U}(n) = i \right]   -   \Pr\left[   \tx{Hyb}_{\mathcal{U}(n)}(1)=1 \mid \mathcal{U}(n) = i \right] \right) \right. \nonumber \\
	&\left. - \sum_{i=0}^{T}  \left(  \Pr\left[   \tx{Hyb}_{\mathcal{U}(n)+1}(0)=1 \mid \mathcal{U}(n) = i \right]   -   \Pr\left[   \tx{Hyb}_{\mathcal{U}(n)}(0)=1 \mid \mathcal{U}(n) = i \right] \right) \right| \nonumber \\
	= & (T+1)\left|  \sum_{i=0}^{T}  \left(  \Pr\left[   \tx{Hyb}_{\mathcal{U}(n)+1}(1)=1 \land \mathcal{U}(n) = i \right]   -   \Pr\left[   \tx{Hyb}_{\mathcal{U}(n)}(1)=1 \land \mathcal{U}(n) = i \right] \right) \right. \nonumber\\ %\label{eq::03}
	&\left. - \sum_{i=0}^{T}  \left(  \Pr\left[   \tx{Hyb}_{\mathcal{U}(n)+1}(0)=1 \land \mathcal{U}(n) = i \right]   -   \Pr\left[   \tx{Hyb}_{\mathcal{U}(n)}(0)=1 \land \mathcal{U}(n) = i \right] \right) \right| \nonumber \\
	&\leq (T+1) \left|  \Pr\left[   \tx{Hyb}_{\mathcal{U}(n)+1}(1)=1\right] - \Pr\left[   \tx{Hyb}_{\mathcal{U}(n)}(1)=1  \right]  \right|  \nonumber \\
	&+  (T+1) \left|  \Pr\left[   \tx{Hyb}_{\mathcal{U}(n)+1}(0)=1\right] - \Pr\left[   \tx{Hyb}_{\mathcal{U}(n)}(0)=1  \right]  \right| \nonumber
\end{align}	
The condition $\tx{Hyb}_{0}(0)  =  \tx{Hyb}_{0}(1)$ trivially holds since all challenge-equal ciphertexts are random ciphertexts. We use this fact in (\ref{eq::02}). 

By the above computation, we have to prove that 
\begin{align*}
	\left|  \Pr\left[   \tx{Hyb}_{\mathcal{U}(n)+1}(b)=1\right] - \Pr\left[   \tx{Hyb}_{\mathcal{U}(n)}(b)=1  \right]  \right| \leq \negl, ~~~\text{where}~b\in\bin.
\end{align*} 


In order to use the firewall technique \cite{jiang2020,nishimaki2022direction}, we define the next hybrid game:
\begin{description}
	\item[$\tx{Hyb}_{i}'(b):$] This is the same as $\tx{Hyb}_{i}(b)$ except that the game chooses $\tx{fwl},\tx{fwr}\gets [0,T]$. If the adversary corrupts $\key_{j}$ such that $j\in[\tx{fwl},\tx{fwr}]$ or $\tok_{\tx{fwr}+1}$, the game aborts.
\end{description}
The guess is correct with probability $(T+1)^{-2}$. Therefore, we have
%
\begin{align*}
	\left|  \Pr\left[   \tx{Hyb}_{\mathcal{U}(n)+1}(b)=1\right]\right. - & \left. \Pr\left[   \tx{Hyb}_{\mathcal{U}(n)}(b)=1  \right]  \right| \\
	&\leq (T+1)^{2}	\left|  \Pr\left[   \tx{Hyb}_{\mathcal{U}(n)+1}'(b)=1\right] - \Pr\left[   \tx{Hyb}_{\mathcal{U}(n)}'(b)=1  \right]  \right|.
\end{align*}
%
Note that to prove the security of our UE scheme it remains to show that the following condition holds.  
\begin{align}\label{eq::04}
	\left|  \Pr\left[   \tx{Hyb}_{\mathcal{U}(n)+1}'(b)=1\right] - \Pr\left[   \tx{Hyb}_{\mathcal{U}(n)}'(b)=1  \right]  \right|\leq \negl.
\end{align}

\begin{lemma}
	If the LWE assumption holds, then (\ref{eq::04}) holds.
\end{lemma}	
%
The proof of this lemma is omitted as it is almost identical to that of \cite[Lemma 5.7]{nishimaki2022direction}.

\vspace{0.5cm}
Our considerations justify the following theorem.
\begin{theorem}
	Under the LWE assumption, the presented $\tx{FrodoPKE}$-based UE scheme is $\tx{rand\md ind\md eu\md cpa}$	secure in the backward-leak uni-directional setting. That is, $\Adv_{\adv}^{\EUS}(n) \leq \negl$. 
\end{theorem}	

%\end{pcvstack}	(\vr{C}_{1}, \vr{C}_{2}) \in  \ZZ_{q}^{\bar{m}\times {n}}  \times \ZZ_{q}^{\bar{m}\times \bar{n}}
\section{Implementations}

\noindent This section provides a detailed analysis of the execution times for different parts of the updatable encryption scheme based on $\tx{FrodoPKE}$ variants, including AES and SHAKE versions for various security levels (i.e. NIST level 1, level 3, and level 5). The results include average times and standard deviations for key generation, encryption, decryption, token generation, and update operations. 

\noindent We made the implementation of an updatable encryption scheme based on $\tx{FrodoPKE}$ in Python programming language utilizing NumPy library. As the core of our implementation, we used the reference one of $\tx{FrodoPKE}$. 
It should be noticed that although Python is not be considered as a high-performance language, the primary goal of this implementation is to compare the execution time of the individual components of an updatable encryption scheme based on $\tx{FrodoPKE}$. Despite Python's slower execution time, it still remains a popular choice for this tasks since its simplicity and the access to a number of libraries. A library that can be indicated here is NumPy, which provides a lot of powerful tools for numerical computations. 

\noindent We conducted 100 series of tests in terms of the execution times,  so as to analyze the relationships between algorithm parts depending on algorithm variants and random messages. Out of the examined factors, decryption times are consistently low and have minimal standard deviations. Due to multiplications of many large matrices, the most computationally complex part of the algorithm is token generation and its computation time increases significantly with the type of $\tx{FrodoPKE}$. These results highlight the trade-offs between security levels and performance, with higher security levels leading to increased computation times, especially in token generation and overall execution time.




\vspace{5mm}
\begin{table}[ht]
	\centering{
		\caption{	\label{TestPhi}  The tests of updatable encryption scheme based on $\tx{FrodoPKE}$ AES. }
		\scalebox{1}{
			\hspace*{-2cm}
			\begin{tabular}{ |p{3.3cm}|p{1cm}|p{1cm}|p{1cm}|p{1cm}|p{1cm}|p{1cm}| }
				\hline
				\multicolumn{1}{|c|}{{Algorithm variants}} &\multicolumn{2}{|c|}{FrodoPKE-640}& \multicolumn{2}{|c|}{FrodoPKE-976}&  \multicolumn{2}{|c|}{FrodoPKE-1344}\\
				\hline
				
				Algorithm parts&average time [s] &std &average time [s] &std &average time [s] &std\\
				
				\specialrule{.2em}{.1em}{.1em} 
				
				UE.KG(params)&   0.0198 & 0.0006   & 0.0465   &0.0030& 0.0878   &0.0035\\
				\hline
				UE.Enc($\text{k}_e, \mathbf{m} \in \mathcal{M}$)&  0.0218   &0.0011   &0.0475& 0.00104& 0.0939   &0.0010\\
				\hline
				UE.Dec($\text{k}_e, \text{ct}_e$)& 0.0003   & 0.00001    & 0.0004   &0.00001 & 0.0004   &0.00001\\
				\hline
				UE.TG($\text{k}_e, \text{k}_{e+1}$)&  6.3251  & 0.2018    & 18.0592   &0.9919&  64.1637  &0.8593\\
				\hline
				UE.Upd($\Delta_{e+1}, \text{ct}_e$)&   0.1224 & 0.0057     & 0.5130  &0.0057& 1.2693   &0.03790\\
				\hline
		\end{tabular}}
		\hspace*{-2cm}
	}
	
	
\end{table}

\begin{table}[ht]
	\centering{
		\caption{	\label{TestPhi2}  The tests of updatable encryption scheme based on $\tx{FrodoPKE}$ SHAKE. }
		\scalebox{1}{
			\hspace*{-2cm}
			\begin{tabular}{ |p{3.3cm}|p{1cm}|p{1cm}|p{1cm}|p{1cm}|p{1cm}|p{1cm}| }
				\hline
				\multicolumn{1}{|c|}{{Algorithm variants}} &\multicolumn{2}{|c|}{FrodoPKE-640}&  \multicolumn{2}{|c|}{FrodoPKE-976}&  \multicolumn{2}{|c|}{FrodoPKE-1344}\\
				\hline
				
				Algorithm parts&average time [s] &std &average time [s] &std &average time [s] &std\\
				
				\specialrule{.2em}{.1em}{.1em} 
				
				UE.KG(params)  &0.0198 &   0.0004  & 0.0471&0.0030&  0.0894   &0.0063\\
				\hline
				UE.Enc($\text{k}_e, \mathbf{m} \in \mathcal{M}$)&     0.0218  & 0.0002     &0.0484& 0.0014& 0.0945   &0.0006\\
				\hline
				UE.Dec($\text{k}_e, \text{ct}_e$) &0.0002 &   0.00001  & 0.0004   &0.00001& 0.0005   &0.00003\\
				\hline
				UE.TG($\text{k}_e, \text{k}_{e+1}$)  &6.3118 &   0.0649 & 18.0746   &0.4328& 63.3840   &0.8250\\
				\hline
				UE.Upd($\Delta_{e+1}, \text{ct}_e$)   &0.1214 &    0.0012  & 0.5166   &0.01533& 1.2572 &0.0267\\
				\hline
		\end{tabular}}
		\hspace*{-2cm}
	}
\end{table}
%\section{Future work}
%One of the most promising directions for future research is the continued utilization of the  $\tx{FrodoPKE}$ scheme.  $\tx{FrodoPKE}$ has proven to be an efficient and highly adaptable tool, particularly for applications that demand high levels of security. Its design, based on lattice-based cryptography, provides robust defenses against quantum attacks, making it a strong candidate for systems requiring post-quantum security. Given this, its inclusion in our work has significantly contributed to the overall effectiveness and security of our approach.
%
%Moreover,  $\tx{FrodoPKE}$'s flexibility makes it suitable for more cryptographic constructions, such as proxy re-encryption. Proxy re-encryption allows an intermediary to transform ciphertexts intended for one party into ciphertexts for another, without accessing the original plaintext. The secure nature of  $\tx{FrodoPKE}$ could potentially make it a strong foundation for building such proxy re-encryption schemes that offer both post-quantum security and efficiency.
%
%A further avenue worth exploring involves adjustments to the probability distributions used in  $\tx{FrodoPKE}$'s core operations and its updateable version.  $\tx{FrodoPKE}$'s security relies heavily on the precise construction of these distributions, particularly in generating error for the Learning with Errors (LWE) problem. 
%
%Overall,  $\tx{FrodoPKE}$ provides a strong foundation for high-security cryptographic applications. However, through targeted improvements—both in its implementation and in exploring additional applications, such as proxy re-encryption—there is significant potential for advancing its capabilities further.

\nocite{*}
\bibliographystyle{fundam}
\bibliography{citations}

%%%%%%%%%%%%%%%%%%%%%%%%%%%%%%%%%%%%%%%%%%%%%%%%%%%%%%%%%%%%%%%%%%%%%%


\end{document}
